% =====================================================================
%  Cac hinh ve lai theo bo ky hieu va bang mau thong nhat
%  Bien dich:  xelatex hinh_thongnhat.tex
%  Gom: Hinh 2.1, Hinh 3.1, Hinh 3.3 (gop 3.3 + 3.4 cu)
% =====================================================================
\documentclass[11pt]{article}
\usepackage[margin=1.4cm,paperwidth=19cm,paperheight=34cm]{geometry}
\usepackage{fontspec}
\setmainfont{Liberation Serif}   % Overleaf: doi thanh "Times New Roman"
\usepackage{amsmath}
\usepackage{amssymb}
\usepackage{tikz}
\usetikzlibrary{positioning,fit,arrows.meta,calc,decorations.pathreplacing,shapes.geometric,backgrounds}
\pagestyle{empty}
\setlength{\parindent}{0pt}

% ---------------------------------------------------------------------
%  BANG MAU DUNG CHUNG CHO MOI HINH TRONG LUAN VAN
%    UE1 cam | UE2 teal | UE3 xanh duong | UE4 tim
%    luong chung magenta | kenh/moi truong green | khoi xu ly gray
% ---------------------------------------------------------------------
\definecolor{cUEa}{RGB}{217,119,6}     % UE1
\definecolor{cUEb}{RGB}{13,148,136}    % UE2
\definecolor{cUEc}{RGB}{37,99,235}     % UE3
\definecolor{cUEd}{RGB}{124,58,237}    % UE4
\definecolor{cChung}{RGB}{190,24,93}   % luong chung
\definecolor{cKenh}{RGB}{22,101,52}    % kenh / STAR-RIS

\tikzset{
  bx/.style  ={draw,fill=gray!12,align=center,inner sep=3pt,font=\scriptsize,line width=0.5pt},
  sm/.style  ={draw,fill=gray!12,align=center,inner sep=2.5pt,font=\tiny,line width=0.45pt},
  ue/.style  ={draw,circle,inner sep=0pt,minimum size=0.5cm,font=\tiny,line width=0.5pt},
  tit/.style ={font=\small\bfseries,anchor=north west},
  note/.style={font=\tiny,align=left,anchor=north west},
}
\newcommand{\chuthich}[2]{\vspace{1.5mm}{\footnotesize #1\\ \textit{Nguồn: #2}}\vspace{9mm}}

\begin{document}
\centering

% =====================================================================
% HINH 2.1 -- mo hinh he thong
% =====================================================================
\begin{tikzpicture}[font=\scriptsize,>=Stealth,line width=0.5pt]

% --- BS mot anten ---
\node[bx,minimum height=0.85cm,minimum width=0.55cm] (bs) at (0.45,1.85) {BS};
\draw[cUEc,line width=0.6pt] (0.72,1.85) -- ++(0.20,0);
\fill[cUEc] (0.92,1.78) -- (0.92,1.92) -- (1.08,1.85) -- cycle;

% --- STAR-RIS ---
\foreach \i in {0,...,10}{
   \fill[cKenh!22] (4.29,0.42+\i*0.28) rectangle ++(0.22,0.22);
   \draw[cKenh,line width=0.35pt] (4.29,0.42+\i*0.28) rectangle ++(0.22,0.22);
}
\draw[cKenh,line width=0.8pt] (4.40,0.34) -- (4.40,3.62);
\node[font=\tiny,anchor=north,align=center] at (4.40,0.26) {STAR--RIS\\ $N$ phần tử};

% --- UE ---
\node[ue,fill=cUEa!18,draw=cUEa] (u1) at (2.25,3.45) {UE$_1$};
\node[ue,fill=cUEc!18,draw=cUEc] (u3) at (2.45,0.55) {UE$_3$};
\node[ue,fill=cUEb!18,draw=cUEb] (u2) at (6.75,3.45) {UE$_2$};
\node[ue,fill=cUEd!18,draw=cUEd] (u4) at (6.95,0.55) {UE$_4$};

% --- duong BS -> STAR-RIS : g ---
\draw[cKenh,line width=0.9pt,->] (1.10,1.85) -- (4.25,2.00);
\node[font=\scriptsize,cKenh,anchor=south] at (2.75,2.02) {$g$};

% --- duong ghep tang STAR-RIS -> UE : h_{r,k} ---
\draw[cUEa,->] (4.32,2.55) -- (2.55,3.28);
\node[font=\tiny,cUEa,anchor=south west] at (3.35,3.00) {$h_{r,1}$};
\draw[cUEc,->] (4.32,1.35) -- (2.72,0.72);
\node[font=\tiny,cUEc,anchor=north west] at (3.35,1.02) {$h_{r,3}$};
\draw[cUEb,->] (4.52,2.55) -- (6.45,3.28);
\node[font=\tiny,cUEb,anchor=south east] at (5.85,3.00) {$h_{r,2}$};
\draw[cUEd,->] (4.52,1.35) -- (6.65,0.72);
\node[font=\tiny,cUEd,anchor=north east] at (5.85,1.02) {$h_{r,4}$};

% --- duong truc tiep BS -> UE : h_{d,k} ---
\draw[gray!60,dashed,->] (1.05,2.10) -- (2.05,3.22);
\node[font=\tiny,gray!45!black,anchor=south] at (1.52,2.98) {$h_{d,1}$};
\draw[gray!60,dashed,->] (1.05,1.45) -- (2.25,0.78);
\node[font=\tiny,gray!45!black,anchor=north] at (1.52,1.02) {$h_{d,3}$};
% toi phia truyen qua: vong qua mep be mat
\draw[gray!60,dashed,->] (1.05,2.25) .. controls (3.0,4.85) and (5.6,4.85) .. (6.60,3.72);
\node[font=\tiny,gray!45!black,anchor=south] at (4.40,4.62) {$h_{d,2}$};
\draw[gray!60,dashed,->] (1.05,1.45) .. controls (3.0,-1.05) and (5.6,-1.05) .. (6.80,0.30);
\node[font=\tiny,gray!45!black,anchor=north] at (4.40,-0.86) {$h_{d,4}$};

% --- nhan hai phia ---
\node[font=\tiny,cKenh,anchor=east] at (4.25,3.82) {phía phản xạ};
\node[font=\tiny,cKenh,anchor=west] at (4.55,3.82) {phía truyền qua};

\node[note,text width=12.5cm] at (0,-1.50)
  {Đường liền: đường ghép tầng BS $\to$ STAR--RIS $\to$ UE.
   Đường nét đứt: đường trực tiếp BS $\to$ UE, tồn tại với mọi $k$.
   Kênh hiệu dụng của UE $k$ là $h^{\mathrm{eff}}_k=h_{d,k}+\sum_{n}h^{*}_{r,k,n}\,\phi^{X}_n\,g_n$
   với $X\in\{T,R\}$ theo phía của UE.};
\end{tikzpicture}

\chuthich{Hình 2.1: Mô hình hệ thống SISO STAR--RIS hỗ trợ RSMA}
{Tác giả xây dựng theo công thức (2.6) (2026)}

\clearpage

% =====================================================================
% HINH 3.1 -- khung DRL (thay the ca Hinh 1.3 cu)
% =====================================================================
\begin{tikzpicture}[font=\scriptsize,>=Stealth,line width=0.5pt]

\node[bx,text width=2.0cm,minimum height=1.0cm] (csi)  at (0,0)    {CSI\\ $h_d,\ g,\ h_r$};
\node[bx,text width=2.2cm,minimum height=1.0cm] (st)   at (3.0,0)  {Tạo trạng thái\\ và chuẩn hóa\\ \tiny $s\in\mathbb{R}^{d_s}$};
\node[bx,text width=2.6cm,minimum height=1.0cm] (net)  at (6.4,0)  {Actor\\ DDPG / TD3 / PPO};
\node[bx,text width=2.4cm,minimum height=1.0cm] (act)  at (9.8,0)  {Hành động chuẩn hóa\\ $a\in[-1,1]^{d_a}$};
\node[bx,fill=cChung!10,text width=2.6cm,minimum height=1.0cm] (dec) at (9.8,-2.3)
      {Bộ giải mã hành động vật lý\\ \tiny $p,\ \eta,\ \beta,\ \theta^T,\ \theta^R$};
\node[bx,fill=cKenh!10,text width=2.8cm,minimum height=1.0cm] (env) at (5.4,-2.3)
      {Môi trường\\ STAR--RIS--RSMA};
\node[bx,text width=2.4cm,minimum height=1.0cm] (rw)  at (1.4,-2.3)
      {Tốc độ, QoS\\ và phần thưởng};

\draw[->] (csi)  -- (st);
\draw[->] (st)   -- (net);
\draw[->] (net)  -- (act);
\draw[->] (act)  -- (dec);
\draw[->] (dec)  -- (env);
\draw[->] (env)  -- (rw);
\draw[->] (rw.north) -- (1.4,-1.15) -- (6.4,-1.15) -- (net.south);
\node[font=\tiny,anchor=south] at (3.9,-1.12) {tín hiệu học (cập nhật trọng số)};

\node[note,text width=12.5cm] at (-0.9,-3.30)
  {Bộ giải mã hành động vật lý \textbf{dùng chung} cho DDPG, TD3 và PPO, nên khác biệt
   giữa ba thuật toán chỉ nằm ở cách cập nhật chính sách, không nằm ở miền hành động.};
\end{tikzpicture}

\chuthich{Hình 3.1: Khung DRL và bộ giải mã hành động vật lý dùng chung}
{Tác giả xây dựng (2026)}

\clearpage

% =====================================================================
% HINH 3.3 -- gop Hinh 3.3 + 3.4 cu thanh mot hinh ba cot
% =====================================================================
\begin{tikzpicture}[font=\scriptsize,>=Stealth,line width=0.5pt]
\def\bw{3.1}   % be rong mot cot
\def\bh{0.52}  % chieu cao mot khoi

% ---------- cot 1: Actor xac dinh (DDPG / TD3) ----------
\begin{scope}[shift={(0,0)}]
  \node[tit] at (0,0.95) {(a) Actor — DDPG, TD3};
  \node[sm,anchor=north,text width=\bw cm] (a1) at (1.55,0.25)  {Trạng thái $s$ \ ($d_s$)};
  \node[sm,anchor=north,text width=\bw cm] (a2) at (1.55,-0.55) {FC 256 \ $\to$ \ LN + ReLU};
  \node[sm,anchor=north,text width=\bw cm] (a3) at (1.55,-1.35) {FC 256 \ $\to$ \ LN + ReLU};
  \node[sm,anchor=north,text width=\bw cm] (a4) at (1.55,-2.15) {FC $d_a$};
  \node[sm,anchor=north,text width=\bw cm,fill=cChung!10] (a5) at (1.55,-2.95) {Tanh \ $\to$ \ $a\in[-1,1]^{d_a}$};
  \foreach \i/\j in {a1/a2,a2/a3,a3/a4,a4/a5}{ \draw[->] (\i.south) -- (\j.north); }
\end{scope}

% ---------- cot 2: Actor ngau nhien (PPO) ----------
\begin{scope}[shift={(4.9,0)}]
  \node[tit] at (0,0.95) {(b) Actor — PPO};
  \node[sm,anchor=north,text width=\bw cm] (b1) at (1.55,0.25)  {Trạng thái $s$ \ ($d_s$)};
  \node[sm,anchor=north,text width=\bw cm] (b2) at (1.55,-0.55) {FC 256 \ $\to$ \ LN + ReLU};
  \node[sm,anchor=north,text width=\bw cm] (b3) at (1.55,-1.35) {FC 256 \ $\to$ \ LN + ReLU};
  \node[sm,anchor=north,text width=\bw cm] (b4) at (1.55,-2.15) {FC $d_a$ \ $\to$ \ $\mu(s)$};
  \node[sm,anchor=north,text width=\bw cm] (b5) at (1.55,-2.95) {Lấy mẫu $u\sim\mathcal{N}(\mu,\sigma^2)$\\ \tiny $\log\sigma$ là tham số học riêng};
  \node[sm,anchor=north,text width=\bw cm,fill=cChung!10] (b7) at (1.55,-3.75) {Tanh \ $\to$ \ $a=\tanh(u)$};
  \foreach \i/\j in {b1/b2,b2/b3,b3/b4,b4/b5,b5/b7}{ \draw[->] (\i.south) -- (\j.north); }
\end{scope}

% ---------- cot 3: Critic ----------
\begin{scope}[shift={(9.8,0)}]
  \node[tit] at (0,0.95) {(c) Critic};
  \node[sm,anchor=north,text width=\bw cm] (c1) at (1.55,0.25)  {$[s,a]$ \ ($d_s+d_a$) — DDPG/TD3\\ $s$ \ ($d_s$) — PPO};
  \node[sm,anchor=north,text width=\bw cm] (c2) at (1.55,-0.80) {FC 256 \ $\to$ \ LN + ReLU};
  \node[sm,anchor=north,text width=\bw cm] (c3) at (1.55,-1.60) {FC 256 \ $\to$ \ LN + ReLU};
  \node[sm,anchor=north,text width=\bw cm] (c4) at (1.55,-2.40) {FC 1};
  \node[sm,anchor=north,text width=\bw cm,fill=cKenh!10] (c5) at (1.55,-3.20)
       {$Q(s,a)$ — DDPG/TD3\\ $V(s)$ — PPO};
  \foreach \i/\j in {c1/c2,c2/c3,c3/c4,c4/c5}{ \draw[->] (\i.south) -- (\j.north); }
\end{scope}

\node[note,text width=12.8cm] at (0,-4.55)
  {Cả ba cột dùng cùng bề rộng ẩn 256 và cùng thứ tự FC $\to$ LN $\to$ ReLU.
   TD3 dùng hai Critic độc lập có cấu trúc giống cột (c) và lấy giá trị nhỏ hơn khi tạo mục tiêu.
   Đầu ra của cả hai Actor đều đi vào cùng một bộ giải mã hành động vật lý ở Hình 3.1.};
\end{tikzpicture}

\chuthich{Hình 3.5: Kiến trúc khối của Actor và Critic}
{Tác giả xây dựng theo cấu hình thực nghiệm (2026)}

\end{document}
